\section{The reciprocal criterion with the physical cutoff fixed}
\label{sec:fixed-cutoff-reciprocal-criterion}

MA-MO-031 varies the cutoff while keeping the test fixed.
Here the cutoff is fixed in advance.  The test varies in
one specified compact-support Frechet space, and the
operator remains the actual fixed-core composition.
No nonpolynomial growth of an arithmetic reciprocal
is asserted.

Retain the generic rational line, finite \(S\),
\(\mathcal H_a,V_a,A,j,\mathscr C\) of (MA-MO.1000).
In particular
\[
 |A^{(k)}(t)|\le C_k(2+|t|)^{D_k},\qquad
 \left|\frac{j'(t)}{j(t)}\right|
                         \le C(2+|t|)^D,\qquad
 |j(t)|\le\exp(C\log^2(2+|t|)).
\tag{MA-MO.1100}
\]
Fix once and for all a real smooth bounded cutoff
\(\phi\), identically one far left and identically
zero far right, with \(\phi'\in C_c^\infty\).
The proof permits \(0\le\phi\le1\), including any
predetermined monotone cutoff of that kind; no
cutoff perturbation will be made.  Put
\[
 T_\phi=\mathcal F_+M_\phi\mathcal F_+^{-1},\qquad
 p(v)=\int e^{ivx}\phi'(x)\,dx,\qquad
 p(0)=-1 .
\tag{MA-MO.1101}
\]
Let \(\mathscr D_K\) be the complex Frechet space
of smooth functions supported in \(K=[1,2]\),
with its usual derivative seminorms, and let
\(\mathscr D_K^\mathbb R\) be its real subspace.
All these tests satisfy \(g(0)=0\).
For \(g\in\mathscr D_K\), write
\[
 m_g(t)=\mathcal Mg(a+it),\qquad
 K^\circ_gv
   =M_{m_g}T_\phi v-M_{m_gA}T_\phi M_jv,
                           \quad v\in\mathscr C .
\tag{MA-MO.1102}
\]
Both terms are literal \(L^2\) functions on this
fixed core, as verified in (MA-MO.1005).

\paragraph{MA-MO-032 (a predetermined cutoff and a fixed test-support space).}
The following are equivalent:
\[
\begin{array}{ll}
\mathrm{(P)}&
 |j(t)|\le C(2+|t|)^N\text{ for some }C,N;\\[2pt]
\mathrm{(B_{\mathbb C})}&
 K^\circ_g\text{ extends boundedly to }L^2
                         \text{ for every }g\in\mathscr D_K;\\[2pt]
\mathrm{(B_{\mathbb R})}&
 K^\circ_g\text{ extends boundedly to }L^2
                         \text{ for every }g\in\mathscr D_K^\mathbb R;\\[2pt]
\mathrm{(T)}&
 K^\circ_g\text{ extends to a trace-class operator}
                         \text{ for every }g\in\mathscr D_K.
\end{array}
\tag{MA-MO.1103}
\]
When these conditions hold, the original sequential
commutator is defined on the whole fixed core for
every compact smooth test, including tests with
other compact supports.  It has the Schwartz kernel
of MA-RS-002, and the actual finite-frequency
operators of MA-RS-007 converge in trace norm.
The operator depends continuously in trace norm
on the test on each fixed compact support.

If (P) fails, the set of \(g\) for which
\(K^\circ_g\) has no bounded extension is a dense
\(G_\delta\) in each of \(\mathscr D_K\) and
\(\mathscr D_K^\mathbb R\).  For every such test
there are unit vectors \(v_n\in\mathscr C\)
with \(\|K^\circ_gv_n\|_2\to\infty\), and the
finite operators for this same fixed \(g,\phi\)
satisfy
\[
                   \sup_N\|K_N(g)\|=\infty ,
\tag{MA-MO.1104}
\]
although their scalar traces converge by MA-RS-007.

If the predetermined \(\phi'\) has finite Gevrey
order, the original sequential composition
\[
                C_gZ_{a,S}[J_{a,S},M_\phi]
\tag{MA-MO.1105}
\]
is defined on \(\mathscr D_a\) for every
\(g\in C_c^\infty\), without assuming (P).
Thus in the conditional failure branch the
unbounded operators are genuine sequential
commutators on their entire fixed core, not
merely alternate formulas used after a failed
inverse-domain condition.

\begin{proof}
\emph{1. A polynomial lower value at a nearby
ordinate, derived from the finite disk.}
Choose \(c>0\), with \(c<1/8\), so small that
\[
 |p(v)|\ge\tfrac12
                      \quad(c/2\le v\le3c).
\tag{MA-MO.1106}
\]
This follows from \(p(0)=-1\) and continuity.
For every sufficiently large \(|T|\), there
is \(v_T\in(c,2c)\) for which, putting
\(t_T=T+v_T\),
\[
                    |A(t_T)|\ge c_0(2+|T|)^{-D_0},
\tag{MA-MO.1107}
\]
with constants independent of \(T\).
We give the full lower-bound argument.

In the raw zeta variable use a disk centered
at \(w_T=2+iT\), with a zero-free boundary
radius \(12<R_T<13\).  Its zero count with
multiplicities is \(M_T=O(\log(2+|T|))\).
For large \(|T|\) it contains neither the
pole nor trivial zeros.  Write
\(\zeta=H_TB_T\), where \(B_T\) is the finite
disk Blaschke product with multiplicities.
The fixed-strip polynomial upper bound and
\(|B_T|=1\) on the boundary give a polynomial
upper bound on \(|H_T|\) throughout the disk.
At its center,
\[
 |H_T(w_T)|\ge|\zeta(2+iT)|\ge\zeta(2)^{-1}.
\]
Harnack's inequality applied to a positive
constant minus \(\log|H_T|\) therefore gives
\[
 |H_T(w)|^{-1}\le C(2+|T|)^C
                       \quad\hbox{on the inner half disk}.
\tag{MA-MO.1108}
\]
This is the actual finite-disk argument of
MA-RS-001, not an assertion of a polynomial
reciprocal along the entire original line.

The points \(w=\sigma+i(T+v)\) with
\(v\in[c,2c]\) lie in the inner quarter.
For a disk zero \(\rho=\beta+i\gamma\),
the normalized factor satisfies
\[
 |B_{\rho,T}(w)|\ge\frac{|w-\rho|}{2R_T}
                 \ge\frac{|T+v-\gamma|}{2R_T}.
\]
Put
\[
 V_T(v)=\sum_{\rho\ {\rm in\ the\ disk}}
                 m_\rho\log^+\frac1{|T+v-\gamma|}.
\tag{MA-MO.1109}
\]
For every real \(\gamma-T\),
\[
 \int_c^{2c}\log^+\frac1{|v-(\gamma-T)|}\,dv
       \le\int_{-1}^1\log(1/|u|)\,du=2 .
\]
Hence \(\int_c^{2c}V_T(v)\,dv\le2M_T\).
There is a point \(v_T\) outside the finitely
many exact ordinates with
\(V_T(v_T)\le(2/c+1)M_T\).
Multiplying the factor estimates yields
\[
 |B_T(\sigma+it_T)|
       \ge(2R_T)^{-M_T}e^{-V_T(v_T)}
                       \ge c(2+|T|)^{-C}.
\tag{MA-MO.1110}
\]
Together with (MA-MO.1108) this gives the
same lower bound for \(|\zeta(\sigma+it_T)|\).
The finitely many factors of \(E_S\) are
bounded and bounded away from zero on the
fixed line \(\sigma>0\), proving
(MA-MO.1107).  No separation among disk
zeros, or simplicity assumption, was used.

\emph{2. Actual separated packets and quantitative
noncancellation.}
Fix a nonnegative \(b\in C_c^\infty((1,2))\)
with \(\int b=1\), and put
\[
 \begin{gathered}
 B(v)=\int e^{ivx}b(x)\,dx,\qquad B(0)=1,\\
 g_T(x)=e^{-ax}e^{-it_Tx}b(x),\qquad
 m_{g_T}(t)=B(t-t_T).
 \end{gathered}
\tag{MA-MO.1111}
\]
These are tests in the one fixed space
\(\mathscr D_K\).  For every integer \(N\),
\[
 \max_{k\le N}\|g_T^{(k)}\|_\infty
                           \le C_N(2+|T|)^N .
\tag{MA-MO.1112}
\]
The off-diagonal distribution kernel of
\(T_\phi\) is the ordinary smooth function
\[
                    \kappa(v)=\frac{i\,p(v)}{2\pi v}
                                      \quad(v\ne0).
\tag{MA-MO.1113}
\]
By (MA-MO.1106), \(\kappa\) is bounded away
from zero, with bounded derivatives, on
\([c/2,3c]\).  Its possible delta component
at zero never enters the separated supports
used below.

There are fixed \(h_0,D_1>0\) such that
\[
                        h_T=h_0(2+|T|)^{-D_1}
\tag{MA-MO.1114}
\]
satisfies all the following for large \(|T|\).
If \(|u-T|\le h_T\) and
\(|t-t_T|\le h_T\), then \(t-u\in[c/2,3c]\)
and each of the four ratios
\[
 B(t-t_T),\qquad
 \frac{A(t)}{A(t_T)},\qquad
 \frac{j(u)}{j(T)},\qquad
 \frac{\kappa(t-u)}{\kappa(v_T)}
\tag{MA-MO.1115}
\]
differs from one by at most \(1/100\).
Here is a quantitative choice argument.
The first and fourth ratios obey the
required estimate on sufficiently small
fixed neighborhoods, by their bounded
derivatives and nonzero lower bounds.
For the second, the polynomial bound for
\(A'\) and (MA-MO.1107) give
\[
 \left|\frac{A(t)-A(t_T)}{A(t_T)}\right|
                         \le C(2+|T|)^{D_A+D_0}h_T .
\]
For the third, integrate \(j'/j\) along
the real segment; since \(j\ne0\),
\[
 \left|\frac{j(u)}{j(T)}-1\right|
       \le \exp(C(2+|T|)^D h_T)-1 .
\]
Choose \(D_1\ge\max\{D_A+D_0,D\}+1\),
and then \(h_0\) small enough, also
less than \(c/8\).  This proves every
claim in (MA-MO.1115).  The product of
the four ratios differs from one by
less than \(1/4\), so its real part
is at least \(3/4\).

Choose a nonnegative \(\alpha\in C_c^\infty((-1,1))\)
with \(\|\alpha\|_2=1\).  Use the Hilbert inner
product linear in its first variable, and set
\[
 \begin{gathered}
 u_T(u)=h_T^{-1/2}\alpha((u-T)/h_T),\\
 C_T=A(t_T)j(T)\kappa(t_T-T),\qquad
 w_T(t)=\frac{C_T}{|C_T|}
               h_T^{-1/2}\alpha((t-t_T)/h_T).
 \end{gathered}
\tag{MA-MO.1116}
\]
All factors in \(C_T\) are nonzero.
The packets are compact smooth, have
norm one, and have separated supports.

For clarity let
\(\mathcal S_T=M_{m_{g_T}A}T_\phi M_j\)
on the core.  The actual separated-support
kernel, the phase in \(w_T\), and
(MA-MO.1115) give
\[
 \begin{split}
 \operatorname{Re}\langle\mathcal S_Tu_T,w_T\rangle
 &\ge\frac34|A(t_T)j(T)\kappa(t_T-T)|\,h_T
                                  \left(\int\alpha\right)^2\\
 &\ge c_1(2+|T|)^{-D_0-D_1}|j(T)|.
 \end{split}
\tag{MA-MO.1117}
\]
This calculation is an ordinary double
integral on two compact separated intervals.
All compositions are defined before that
calculation.  It proves noncancellation by
a common phase and positive packet profiles;
it does not estimate an unsigned kernel
and assume that its operator norm is large.

The other term in (MA-MO.1102) has norm
at most \(\|B\|_\infty\|\phi\|_\infty
\le\|\phi\|_\infty\), independent of \(T\).
Thus, whenever \(K^\circ_{g_T}\) has a
bounded extension, and also as a bound
on its core operator norm if not,
\[
 \|K^\circ_{g_T}\|
       \ge c_1(2+|T|)^{-D_0-D_1}|j(T)|
                              -\|\phi\|_\infty .
\tag{MA-MO.1118}
\]

\emph{3. The universal fixed-support condition
has a finite derivative estimate.}
Assume \(\mathrm{(B_{\mathbb C})}\), and denote
the unique bounded extensions by \(\mathcal K(g)\).
The map
\(\mathcal K:\mathscr D_K\to\mathcal B(L^2)\)
is linear and has closed graph.
Indeed, suppose \(g_n\to g\) in \(\mathscr D_K\)
and \(\mathcal K(g_n)\to T\) in operator norm.
The Mellin transforms converge in every
Schwartz seminorm.  For each core vector
\(v\), both \(T_\phi v\) and \(T_\phi jv\)
are Schwartz, and \(A\) is a polynomial
Schwartz multiplier.  Consequently
\[
           K^\circ_{g_n}v\longrightarrow K^\circ_gv
                                      \quad\hbox{in }L^2.
\tag{MA-MO.1119}
\]
The operator-norm convergence gives
\(Tv=K^\circ_gv\).  Density and uniqueness
of the bounded extension imply
\(T=\mathcal K(g)\), proving closedness.

The finite-seminorm consequence can also
be proved directly, to make its uniform
content explicit.  The sets
\[
 F_n=\{g\in\mathscr D_K:
        \|K^\circ_gv\|_2\le n\|v\|_2
                              \text{ for every }v\in\mathscr C\}
\tag{MA-MO.1120}
\]
are closed, by (MA-MO.1119), and cover
\(\mathscr D_K\) under the assumption.
Baire's theorem gives interior to one
\(F_n\).  Differences give a neighborhood
of zero with operator norm at most \(2n\).
Homogeneity then yields \(C,N\) with
\[
              \|\mathcal K(g)\|
                  \le C\max_{k\le N}\|g^{(k)}\|_\infty
                                      \quad(g\in\mathscr D_K).
\tag{MA-MO.1121}
\]
Thus no uniform estimate over tests has
been assumed without proof.

Apply (MA-MO.1121) to \(g_T\) and use
(MA-MO.1112), (MA-MO.1118).  One obtains
\[
                    |j(T)|\le C'(2+|T|)^{N+D_0+D_1}
\tag{MA-MO.1122}
\]
for large \(|T|\).  The compact range is
bounded by smooth nonvanishing on the line.
This proves (P).

If only \(\mathrm{(B_{\mathbb R})}\) is
assumed, write each complex test as the
sum of its real part and \(i\) times its
imaginary part.  Core linearity gives a
bounded extension for the sum.  Hence
\(\mathrm{(B_{\mathbb C})}\) follows, and
the complex modulated tests in
(MA-MO.1111) are now legitimate.  This
does not silently assume complex test
bounds from a real assertion.

Conversely, (P) gives polynomial growth
of every derivative of \(j\), as proved
in MA-MO-031.  The actual sequential
domain and Schwartz-kernel proof there
applies to this fixed cutoff and every
compact smooth test.  It proves (T),
the trace-norm finite-frequency convergence
and test continuity.  Since (T) implies
\(\mathrm{(B_{\mathbb C})}\), all
conditions in (MA-MO.1103) are equivalent.

\emph{4. Residual genuine test failures
when the reciprocal bound fails.}
Suppose (P) fails.  Define the closed
sets \(F_n\) in (MA-MO.1120) without
assuming that they cover the test space.
Their union is exactly the set of tests
whose core operator has a bounded
extension: the core is dense, and a
uniform bound on it gives that unique
extension.  No \(F_n\) has interior.
Otherwise differences and scaling would
make every test have a bounded extension,
contrary to step3.  Thus their complement
is a dense \(G_\delta\).  The identical
argument on the real Frechet space works,
because the resulting universal real
bound would complexify as above.

For every test in this complement,
unboundedness on the fixed core supplies
unit \(v_n\in\mathscr C\) with
\(\|K^\circ_gv_n\|_2\ge n\).
These are actual core vectors and actual
outputs of (MA-MO.1102).  The choice of
the fixed bad test is by Baire's theorem;
no explicit numerical arithmetic example
is claimed.  Each test is supported in
\([1,2]\), so its value at zero vanishes.

For this same fixed \(g,\phi\), the two
bounded compositions of MA-RS-007 obey
\[
 V_aK_N(g)V_a^{-1}v
              =\chi_NK^\circ_gv
                       \longrightarrow K^\circ_gv
                       \quad(v\in\mathscr C)
\tag{MA-MO.1123}
\]
once \(N\) is large enough for its input
cutoff to equal one on the compact support.
Uniformly bounded operator norms would
therefore bound the core operator, a
contradiction.  This proves (MA-MO.1104)
and excludes strong or trace-norm
convergence to a bounded operator.
The polynomial logarithmic derivative
still gives the actual convergent trace
formula (MA-RS.56), with its residue
correction unchanged.

\emph{5. A fixed Gevrey cutoff keeps the
original sequential domain for every smooth test.}
This last assertion does not use (P).
For \(v\in\mathscr C\), both \(v\) and
\(jv\) are compact smooth.  Outside their
fixed frequency supports, the formula
(MA-RS.20) expresses \(T_\phi v\) and
\(T_\phi jv\) using \(p(t-u)/(t-u)\).
Gevrey decay of \(p\) dominates the
\(\exp(C\log^2(2+|t|))\) growth of \(j\).
Thus \(M_jT_\phi v\in L^2\), and
\(M_AM_jT_\phi v=T_\phi v\in L^2\).
Also \(T_\phi jv\) is in the maximal
domain of \(M_A\).  This checks both
sequential compositions on the entire
core.  The left convolution \(C_g\)
is bounded on \(\mathcal H_a\) for
every compact smooth \(g\), with norm
at most \(\int|g(x)|e^{ax}\,dx\).
No Gevrey estimate on \(g\) enters these
domain checks.  Consequently (MA-MO.1105)
is genuinely defined and equals
\(V_a^{-1}K^\circ_gV_a\) on \(\mathscr D_a\).
Its conditional unboundedness follows
from the preceding actual core witnesses.
\end{proof}

\paragraph{Arithmetic inputs and exact remaining question.}
The stipulated rational input \(A\) places
all nontrivial zeros on the critical line.
For every \(\sigma\in(1/2,1)\) considered
here, it therefore supplies a fixed
horizontal gap and proves (P) by
MA-RS-004.  Under \(A\), every positive
conclusion above is available.  The
failure branch is an explicit conditional
operator theorem; it neither asserts
nonpolynomial reciprocal growth nor a
zero configuration causing it.  The
sufficient multiplicity/distance condition
of (MA-MO.1023) remains conditional too.

Compared with MA-MO-031, the new result
settles the quantifier in which the
physical cutoff is prescribed and the
compact test varies.  It still does
not decide boundedness for one fixed
individual pair \(g,\phi\), or prove
the reciprocal bound independently of
zero-location input.  It is not an
unconditional arithmetic resolution
of RF002/RS003 or FT002.  Singular
boundary distributions and the actual
two-sided ideal return have their
separately proved domains; no packet
above is asserted to belong to that
completed ideal.

The finite-disk argument was derived
in (MA-MO.1108)--(MA-MO.1110) from the
actual polynomial strip bounds, Euler
lower bound and local zero count used
in the supplied MT-M18.7--9 and
MA-RS-001.  The normalized off-diagonal
kernel is (MA-RS.20)/(MA-RS.53).
The core and finite operators are
exactly (MA-RS.8--9)/(MA-RS.54);
the polynomial logarithmic derivative
is (MA-RS.64).  No whole-line
polynomial reciprocal estimate is
hidden in the selected nearby ordinate.

In the inherited cote~1 BR-TR-01/02
framework, scans~079--080 and~111--112,
the distinction remains one between
actual defined continuous/nuclear
maps and a scalar trace distribution.
The new fixed-cutoff criterion concerns
that analytic structure.  It changes
no finite-monodromy representation,
Euler determinant or motivic weight,
and supplies no positivity assertion
from convergence of scalar traces.
These archive locators are inherited.
